\section{Two classical estimates for the quartic equation}
\label{sec:classical-estimates}

Both the geometric deduction and the analytical classification use the
same gradient bound and energy monotonicity. We record their short
proofs here, independently of either rigidity theorem. Let $n=m+1\ge2$
and let $v\in C^2(\R^n,[-1,1])$ solve $\Delta v=v^3-v$.
Throughout,
\[
 W(t)=\frac14(1-t^2)^2,\qquad s_0=\frac{2\sqrt2}{3},
\]
and $\omega_m$ denotes the volume of the unit ball in $\R^m$.
Interior elliptic estimates bootstrap $v$ to a smooth function with
uniform bounds for every derivative of fixed order.

\subsection{Modica's gradient estimate}

We use Modica's gradient bound \cite{Modica1985},
\begin{equation}\label{eq:diffuse-modica}
 \frac12|\nabla v|^2\le W(v).
\end{equation}
For completeness, its entire-solution argument in this normalization is
short. For $P=|\nabla v|^2/2-W(v)$, the identities
$\nabla P=\nabla^2v\nabla v-W'(v)\nabla v$ and
$\Delta P=|\nabla^2v|^2-W'(v)^2$ imply, wherever $\nabla v\ne0$,
\[
 \Delta P\ge
 \frac{2W'(v)}{|\nabla v|^2}\nabla v\cdot\nabla P.
\]
If $\sup P>0$, translate a maximizing sequence and pass to a locally smooth
limit. Its $P$ attains a positive maximum. The strong maximum principle
makes $P$ constant on the component of $\{\nabla v\ne0\}$ containing that
point. This component is also closed: on it
$|\nabla v|^2\ge2\sup P>0$. It is therefore all of $\R^n$. The globally
bounded smooth vector field $\nabla v$ has complete forward flow, and along
that flow $v$ increases at rate at least $2\sup P$, contradicting its
boundedness. This proves \eqref{eq:diffuse-modica}.

\subsection{Normalized energy and monotonicity}

We recall Modica's energy monotonicity formula \cite{Modica1989} and
give the stress-tensor calculation in the present normalization. Set
\begin{equation}\label{eq:finite-density}
 e(v)=\frac12|\nabla v|^2+W(v),\qquad
 M_r(y)=\frac{1}{s_0\omega_m r^m}\int_{B_r(y)}e(v),\qquad
 M_\infty=\sup_{r>0}M_r(0).
\end{equation}
The stress tensor
\(e(v)\Id-\nabla v\otimes\nabla v\) is divergence free.  Its contraction with
the derivative of \(x-y\), followed by the divergence theorem, gives
\[
 r\int_{\partial B_r(y)}\bigl(e(v)-|\partial_rv|^2\bigr)
   =\int_{B_r(y)}\bigl((m+1)e(v)-|\nabla v|^2\bigr).
\]
Consequently
\begin{equation}\label{eq:finite-monotonicity}
 M_r'(y)=\frac{1}{s_0\omega_m}
 \left\{r^{-m-1}\int_{B_r(y)}\left(W(v)-\frac12|\nabla v|^2\right)
       +r^{-m}\int_{\partial B_r(y)}|\partial_rv|^2\right\}\ge0.
\end{equation}
The formula initially holds for almost every \(r\) and hence in integrated
form for every pair of radii.  Containment of shifted balls shows that the
limit at infinity is independent of \(y\), and \(M_r(y)\le M_\infty\).

The derivation uses neither stability nor a density-growth hypothesis.

